% !TeX root = ../navier-stokes-manuscript.tex
\subsection{The spatial ascent inside the Tower}\label{sub:ThePressureCompleteMixedDerivativeTower}

The mixed recurrence gives the argument two independent directions.
Critical scaling shows where they line up.
For
\[
  E_s(t):=\frac12\norm{\Lambda^su(t)}_2^2,
  \qquad
  E_{s,\lambda}(t):=\frac12\norm{\Lambda^su_\lambda(t)}_2^2,
  \qquad
  H=E_{1/2},
\]
the Sobolev calculation from the stretching vortex gives
\begin{equation}\label{eq:BodySpatialEnergyScaling}
  E_{s,\lambda}(t)
  =\lambda^{2s-1}E_s(\lambda^2t).
\end{equation}
At \(s=\tfrac12\), the exponent \(2s-1\) vanishes and the scale factor equals one:
\[
  H_\lambda(t):=E_{1/2,\lambda}(t)=H(\lambda^2t).
\]
Differentiating this identity \(n\) times in \(t\) gives
\begin{equation}\label{eq:BodyCriticalHeightTemporalScaling}
  H_\lambda^{(n)}(t)
  =\lambda^{2n}H^{(n)}(\lambda^2t).
\end{equation}
Now set \(s=n+\tfrac12\) in \eqref{eq:BodySpatialEnergyScaling}:
\begin{equation}\label{eq:BodySpatialHeightScalingDiagonal}
  E_{n+1/2,\lambda}(t)
  =\lambda^{2n}E_{n+1/2}(\lambda^2t).
\end{equation}
The \(n\)-th temporal derivative of critical height and the spatial energy at height \(n+\tfrac12\) have the same scale weight \(\lambda^{2n}\).
The heat clock already contained this match: one time derivative costs \(\lambda^2\), and one additional Sobolev derivative contributes \(\lambda^2\) to the quadratic energy.

The temporal rows enter the height derivatives through an exact quadratic contraction.
Since
\[
  H(t)=\frac12\pair{\Lambda^{1/2}V_0(t)}{\Lambda^{1/2}V_0(t)}_{L^2},
\]
Leibniz's rule gives
\begin{equation}\label{eq:BodyHeightDerivativeTemporalContraction}
  H^{(n)}(t)
  =\frac12\sum_{a=0}^{n}\binom{n}{a}
  \pair{\Lambda^{1/2}V_a(t)}
       {\Lambda^{1/2}V_{n-a}(t)}_{L^2}.
\end{equation}
Thus \(H^{(n)}\) is built from the first \(n\) Eulerian response rows.
It remains a scalar quadratic functional of those rows, distinct from \(V_n\).
This is the bridge that lets the temporal Tower be read at critical height.

Navier--Stokes places the higher spatial energy on the other side of that bridge.
Apply \(\Lambda^s\) to the momentum equation and pair with \(\Lambda^su\).
The time derivative gives \(E_s'\), while viscosity gives
\[
  \nu\pair{\Lambda^su}{\Lambda^s\Delta u}_{L^2}
  =-\nu\norm{\Lambda^{s+1}u}_2^2
  =-2\nu E_{s+1}.
\]
Incompressibility cancels the global pressure pairing after integration by parts.
Retain the complete transport--pressure term as
\[
  \mathcal P_s
  :=-
  \operatorname{Re}
  \pair{\Lambda^su}
  {\Lambda^s\bigl((u\cdot\nabla)u+\nabla p\bigr)}_{L^2}.
\]
The exact energy row is
\begin{equation}\label{eq:BodyCriticalEnergyBalancePreview}
  E_s'=\mathcal P_s-2\nu E_{s+1}.
\end{equation}
At the critical height, one temporal derivative exposes the next spatial height:
\[
  H'=\mathcal P_{1/2}-2\nu E_{3/2}.
\]
Differentiate once more and use \eqref{eq:BodyCriticalEnergyBalancePreview} at \(s=\tfrac32\):
\begin{align*}
  H''
  &=\partial_t\mathcal P_{1/2}-2\nu E_{3/2}'
  \\
  &=\partial_t\mathcal P_{1/2}
  -2\nu\mathcal P_{3/2}
  +(2\nu)^2E_{5/2}.
\end{align*}
The next derivative uses the row at \(s=\tfrac52\):
\[
  H'''
  =\partial_t^2\mathcal P_{1/2}
  -2\nu\partial_t\mathcal P_{3/2}
  +(2\nu)^2\mathcal P_{5/2}
  -(2\nu)^3E_{7/2}.
\]
Every earlier transfer remains in the next row, and each differentiation adds one higher viscous energy.
Induction gives, for \(n\geq1\),
\begin{equation}\label{eq:BodyCriticalHeightSpatialAscentPreview}
  H^{(n)}
  =(-2\nu)^nE_{n+1/2}
  +\sum_{j=0}^{n-1}
  (-2\nu)^j
  \partial_t^{\,n-1-j}\mathcal P_{j+1/2}.
\end{equation}
Because \(\nu>0\), the coefficient of the top spatial energy is nonzero at every order.
The same identity can be read as
\[
  E_{n+1/2}
  =(-2\nu)^{-n}
  \left(
    H^{(n)}
    -\sum_{j=0}^{n-1}
      (-2\nu)^j
      \partial_t^{\,n-1-j}\mathcal P_{j+1/2}
  \right).
\]
Controlling the critical response and every retained transfer term at depth \(n\) controls the spatial energy at height \(n+\tfrac12\).
The identities locate that control; the rectangle estimate in the next section supplies it.

\Needspace{18\baselineskip}
\begin{center}
\captionsetup{hypcap=false}
\captionof{table}{The temporal response and the spatial height exposed at each order.  The final two columns share the scaling factor \(\lambda^{2n}\), while \(V_n\), \(H^{(n)}\), and \(E_{n+1/2}\) remain distinct objects.}
\label{tab:BodyTemporalSpatialDiscovery}
\small
\renewcommand{\arraystretch}{1.4}
\begin{tabularx}{\textwidth}{@{}c c c c@{}}
\toprule
order & Eulerian response & critical response & spatial height exposed \\
\midrule
\(0\) & \(V_0=u\) & \(H=E_{1/2}\) & \(E_{1/2}\) \\
\(1\) & \(V_1=u_t\) & \(H'\) & \(E_{3/2}\) \\
\(2\) & \(V_2=u_{tt}\) & \(H''\) & \(E_{5/2}\) \\
\(3\) & \(V_3=u_{ttt}\) & \(H'''\) & \(E_{7/2}\) \\
\(n\) & \(V_n=\partial_t^nu\) & \(H^{(n)}\) & \(E_{n+1/2}\) \\
\bottomrule
\end{tabularx}
\end{center}

This is where the reversal appears.
The singularity route asks the velocity field to make increasingly sharp changes across increasingly short intervals.
The temporal Tower measures those responses.
At the same order, viscosity places a still higher spatial energy in the completed height row.
Once the finite response and transfer terms are controlled, moving deeper in time produces control farther across the Sobolev scale.
The result is increasing spatial regularity at increasing temporal depth.
The singularity pressure and the viscous ascent are balanced and opposing: one sharpens toward loss of regularity, while the other reaches higher Sobolev spaces.

Each completed row reaches one finite spatial height.
All finite heights at one temporal order form the rowwise target
\[
  V_r\in\bigcap_{m=0}^{\infty}L^2(I;H^m(\R^3)),
\]
on a common time interval \(I\).
Requiring the same at every finite temporal depth gives
\begin{equation}\label{eq:BodyDesiredMixedIntersection}
  u\in
  \bigcap_{r=0}^{\infty}
  \bigcap_{m=0}^{\infty}
  H^r(I;H^m(\R^3)).
\end{equation}
This intersection asks for every finite coordinate while keeping the temporal and spatial orders independent.
Each row carries one finite height at a time; the all-orders regularity belongs to their intersection.

The geometry is a grid:
\[
\begin{array}{c|ccccc}
 & H^0_x & H^1_x & H^2_x & H^3_x & \cdots\\
\hline
V_0=u      & \bullet & \bullet & \bullet & \bullet & \cdots\\
V_1=u_t    & \bullet & \bullet & \bullet & \bullet & \cdots\\
V_2=u_{tt} & \bullet & \bullet & \bullet & \bullet & \cdots\\
V_3=u_{ttt}& \bullet & \bullet & \bullet & \bullet & \cdots\\
\vdots     & \vdots  & \vdots  & \vdots  & \vdots  & \ddots
\end{array}
\]
Moving down selects a later response; moving across asks for a higher Sobolev height inside that response.
Fixing a temporal depth and a spatial height cuts out a finite rectangle.
Larger cutoffs retain every coordinate in the smaller block and add the next rows or columns.

The initial datum supplies the lower edge of every such finite block.
Because \(u_0\in\Ssigma\), \(V_0(0)\) belongs to every finite Sobolev space.
Assume the velocity rows through \(V_n(0)\) have this property and fix a target order \(m\).
Sobolev multiplication at a sufficiently high finite order and the Riesz-transform formula \eqref{eq:BodyTemporalPressureRecurrence} give \(Q_n(0)\in H^{m+1}\).
The recurrence \eqref{eq:BodyTemporalVelocityRecurrence} then uses \(V_n(0)\in H^{m+2}\), the finite transport products, and \(\nabla Q_n(0)\in H^m\) to place \(V_{n+1}(0)\) in \(H^m\).
Induction yields
\begin{equation}\label{eq:BodyInitialPressureCompleteMixedJet}
  J_{n,\alpha}(0)\in H_\sigma^m(\R^3),
  \qquad
  \Pi_{n,\alpha}(0)\in H^m(\R^3)
\end{equation}
for every finite \(n\), \(\alpha\), and \(m\).
This gives the finite coordinates at the initial edge.

What remains is their common control over the interval.
The scaling diagonal and completed height rows have shown which coordinates matter; the established datum-generated theorem now supplies their finite rectangle bounds and their compatible assembly.
The next section turns that theorem into the spatial regularity predicted by the reversal.
